\documentclass[11pt,a4paper]{article}
\usepackage[utf8]{inputenc}
\usepackage[T1]{fontenc}
\usepackage{lmodern}
\usepackage{microtype}
\usepackage[margin=1in]{geometry}
\usepackage{amsmath,amssymb,amsthm}
\usepackage{booktabs}
\usepackage{array}
\usepackage{enumitem}
\usepackage{xcolor}
\usepackage[colorlinks=true,linkcolor=blue!50!black,citecolor=blue!50!black,urlcolor=blue!50!black]{hyperref}

\theoremstyle{plain}
\newtheorem{theorem}{Theorem}
\newtheorem*{theorem2p}{Theorem 2$'$}
\newtheorem*{theorem4p}{Theorem 4$'$}
\theoremstyle{definition}
\newtheorem*{axiom}{Axiom}
\theoremstyle{remark}
\newtheorem*{remark}{Remark}

\newcommand{\tr}{\operatorname{tr}}
\newcommand{\E}{\mathbb{E}}
\newcommand{\R}{\mathbb{R}}
\newcommand{\C}{\mathbb{C}}
\newcommand{\RP}{\mathbb{RP}}
\newcommand{\CP}{\mathbb{CP}}
\newcommand{\SO}{\mathrm{SO}}
\newcommand{\SU}{\mathrm{SU}}
\newcommand{\PU}{\mathrm{PU}}
\newcommand{\PSp}{\mathrm{PSp}}
\newcommand{\PSO}{\mathrm{PSO}}
\newcommand{\Sp}{\mathrm{Sp}}
\newcommand{\U}{\mathrm{U}}
\newcommand{\adj}{\mathrm{adj}}
\newcommand{\Sym}{\mathrm{Sym}}
\newcommand{\diag}{\operatorname{diag}}

\title{An Operational Representation of Systems:\\ A Worked Construction in Lattice Gauge Theory}
\author{Edward Chalk\\ \small Sapientronic \quad \href{mailto:edward.chalk@sapientronic.ai}{edward.chalk@sapientronic.ai}}
\date{18 October 2026}

\begin{document}
\maketitle

\begin{abstract}
A representation of systems in general---not of fields, not of particles---is carried into lattice gauge theory and tested there. The representation, which we call MINOS (a \emph{minimal ordinal systems representation}), begins from the smallest account of a system we can state. Something drives it; something constrains it; and how the two cohere into conduct is supplied by the system itself. Its primitive is a binary comparison. Its rule of evidence is that only the shares of a state can be read; its frame and its phase cannot. From that rule, with every assumption stated, the structures of lattice gauge theory are produced rather than named. The relation between two elements' frames is a group element, and the group is classified. The unreadability of frames is a local gauge freedom, and Gauss's law. The readable loss of carrying a state round a closed path, averaged over the state, is a Wilson character action: in the adjoint representation when the amplitudes are complex, in the traceless symmetric representation when they are real. Under the stated loss rule, therefore, the readable state geometry determines the representation content of the action. The invariant information metric on relations has the Casimir operator as its Laplacian. If the readable poles are kept live rather than averaged, they are projective matter, and the same construction fixes a definite higher-order interaction between plaquettes that share an element. Its coefficients and angular dependence follow from the model and are not fitted, and it is verified by direct diagonalisation. The paper's purpose is methodological and comprehensional. Methodological: a representation devised for systems can be stated precisely, carried into an established mathematical theory, and used there to construct and analyse nontrivial structure without renaming the theory's objects. Comprehensional: a difficult theory's structures become intelligible as consequences of a few statements about systems in general. What the paper does not do is also stated. It does not derive the continuum theory. It says nothing about the mass gap. It leaves complex amplitudes and the composition of independently specified systems as open hypotheses. And it claims no novelty for the representation theory it uses.
\end{abstract}

\noindent\small Every computation referred to is reproduced by the accompanying code (\texttt{minos\_yang\_mills.py}, with its output beside it), which travels with the paper. \copyright\ 2026 Edward Chalk. All rights reserved.
\normalsize

\section{The problem, and the purpose of the paper}

Systems thinking has a vocabulary and physics has mathematics, and the two rarely meet on terms the second would accept. The question this paper asks is narrow and testable. \textbf{Can a representation built to describe systems in general be carried into a demanding domain of mathematical physics? Can it produce structure there that is definite and calculable? And can it withstand the scrutiny that domain applies: stated assumptions, proofs, and checks that can fail?} The domain chosen is lattice gauge theory. Its structures are exact, its objects are finite-dimensional, and its standards are exacting.

The purpose is therefore twofold, and neither part is a claim about nature. It is \textbf{methodological}. The paper is a test of a representation, not a derivation of a physical theory. What it establishes is what the representation can be made to do: be stated as axioms, be applied without renaming, produce consequences that can be computed and checked, and be corrected where it is wrong. And it is \textbf{comprehensional}. The paper offers a way of understanding the structures of lattice gauge theory---gauge covariance, Gauss's law, the choice of representation in the action, the Casimir kinetic term---as consequences of a few statements about systems in general. A reader who understands systems can then see why the theory has the shape it has. A reader who understands the theory can see what a systems account adds and what it merely restates. In short, the aim is that \textbf{the systems representation becomes mathematically respectable}: stated in a form mathematics can work with, and shown to hold up when it does. Nothing here is a prediction. Nothing here is offered as the unique or necessary representation of systems. It is one proposed representation, made precise, and put to a hard test.

The answer we report is positive and qualified, and the qualification is the method. Every step from the representation to a lattice structure either follows from the stated axioms or rests on a hypothesis we name; the hypotheses are few and are listed in \S\ref{sec:axioms}. The results are of three kinds: what is derived, what is reconstructed under assumptions, and what is a consequence of a specified model. We keep the three apart throughout, and \S\ref{sec:interp} says, for each result, which it is.

Three results carry the argument. \textbf{First,} under a stated return-loss rule, the readable state geometry determines the representation content of the lattice action, and determines it differently for real and for complex amplitudes (\S\ref{sec:central}). \textbf{Second,} the same rule, with the relativity of frames, gives the gauge geometry: the group, the local freedom, Gauss's law, and the Casimir kinetic term through the invariant information metric (\S\S\ref{sec:geometry}--\ref{sec:central}). \textbf{Third,} the pure lattice gauge theory the construction yields is an \emph{averaged model}, in which the loop's loss is defined as its average over the readable pole. Equivalently, it is the leading term of a limit in which the poles' own stiffness grows. The full object is a gauge theory coupled to projective matter, the poles themselves. With the poles kept as dynamical variables, the same construction fixes a particular interaction between neighbouring plaquettes whose form is not a free choice (\S\ref{sec:consequence}). The first result shows that the representation makes a distinction with consequences for the action. The second shows that it organises known structure under stated assumptions. The third shows that, once a model is specified, it determines quantities that can be calculated and checked. Everything is checked (\S\ref{sec:verification}), and the checks are supplied.

What the paper does not claim is as important. It does not derive Yang--Mills theory from the bare notion of a system. Complex amplitudes are a hypothesis; so is the kinetic-plus-potential form of the dynamics; so is the loss functional. It makes no statement about the continuum limit or the mass gap. It does not claim that the representation theory it uses---character decompositions, invariant integration, degenerate perturbation theory---is new. And it leaves open the general composition law for independently specified systems, which is where the representation's own content would next be tested.

\section{The primitive}

\subsection{The triad}

Consider what any account of a system must contain. Something \textbf{drives} it: a flow, a pressure, a demand, a field. Something \textbf{constrains} it: a bound, a rule, a conservation law, a capacity. And the two do not meet of themselves: how a given drive and a given constraint cohere into conduct is contained in neither and must be supplied---we call it the \textbf{relation}. The three answer three different questions: what moves it, what holds it, how the two cohere. Each is empty without the other two. A drive alone is a tendency with nowhere to go; a constraint alone is a boundary with nothing pressing on it; a relation alone relates nothing.

Three and not two, because a drive and a constraint do not determine what happens; the same pressure against the same wall gives a leak, a burst or a standing load, and what decides is a part of the system, not of its surroundings. Three and not four, because anything further one might name---a state, an output, a history---is a reading of these three or a composite of such triads. The claim of \emph{minimality} is this: the triad is the smallest set of roles under which the word \emph{system} applies, and the representation builds everything from the roles and nothing from their contents.\footnote{\emph{A note on the name.} The acronym was chosen for its words. The fit with the Minos of classical mythology is close enough to record and too slight to lean on. The king is remembered for his \textbf{law}, received and administered, and for his \textbf{royal will} and what its excesses cost him. The third thing remembered of him, that he judges the passage of the dead, is not a third attribute. It is what the meeting of the first two made of him: by his law a judge, by his will a man who had himself made the passage. The triad has the same form: a constraint, a drive, and a relation that is not a third ingredient but what their meeting supplies. Nothing in the paper rests on it.}

Two features of the triad are used throughout. \textbf{The roles are relative, not intrinsic.} What is drive to one element is constraint to the one before it---the pressure that drives a valve is the pressure the pump is constrained to supply---so the names belong to the reading, not to the thing. \textbf{The relation is the system's own contribution.} Being supplied in neither pole, it is the one part that can be anything consistent with the poles. That is why, once the poles are given as a frame, the admissible relations form a group acting on that frame, and why a chain of elements carries a connection along it.

\subsection{The element as a binary comparison}\label{sec:comparison}

The simplest thing a system can do with a drive and a constraint is apportion a unit between them: so much passed on the drive's side, so much held on the constraint's side. If the two shares sum to a whole, the apportioning is one angle: passed $=\cos^2\theta$, held $=\sin^2\theta$. An element, in this representation, \emph{is} that angle. Its two ends are the two ways a system fails: all drive and no constraint, which runs on without bound, and all constraint and no drive, which holds and passes nothing. Between them lies every system that works.

Written out, the element is a \textbf{binary comparison} of a \emph{state} against a \emph{measurement}: with $Q_c=|c\rangle\langle c|$ the state and $P_d=|d\rangle\langle d|$ the passing outcome,
\begin{equation}
\text{passed } p=\tr(P_dQ_c)=|\langle d,c\rangle|^2,\qquad \text{held } h=\tr\big((I-P_d)Q_c\big)=1-p.
\end{equation}
The two are the exclusive outcomes of one binary measurement on one state, which is why $p+h=1$ whatever the angle between $d$ and $c$, and why $d=c$ gives $p=1$ without $c$ having gone anywhere. For a space of more than two directions the held outcome $I-P_d$ is a subspace, not a direction; the two-pole element is the case in which it is one. A system with several interfaces is not a primitive element but a \textbf{composite of comparisons}, one per interface, each with its own normalisation; this matters in \S\ref{sec:geometry} and is why several pass shares at one site may sum to more than one.

\subsection{What is readable, and the relativity of frames}

The representation's one rule of evidence is that \textbf{only shares are readable}. The oriented angle $\theta$, the amplitudes $(\cos\theta,\sin\theta)$ with their signs, the frame in which they are written and any phase they carry are not directly readable. The unoriented acute angle $\arccos\sqrt p$ is, because the share determines it, and so is anything else the shares determine. Two states that give the same shares against every measurement are one readable state. Two consequences make the representation \emph{ordinal} rather than metric. First, $\theta$ is an angle in a plane that belongs to the element. There is no common frame in which two elements' planes could be compared directly, so the relation between two elements' frames is itself something to be supplied, across every adjacency. We call this \textbf{the relativity of frames}. Second, nothing in the element carries a magnitude: a share is a share of a unit, and the unit is the element's own. If a system is to have a size, the size must be made by its composition. In a theory whose only scale is self-generated, a representation that imports none is the right starting point.

\section{The axioms, and what is assumed}\label{sec:axioms}

We state the representation as axioms once, in the form the results below use. Each is marked either as a \emph{postulate} of the proposed representation or as an \emph{assumption} added for the domain. A postulate is what we propose a system is, mathematically; it is not a consequence of the everyday meanings of drive and constraint, and is not claimed as one. The hypotheses are listed so that a reader can see every place where the chain from systems to gauge theory is conditional.

\begin{description}[leftmargin=1.5em,style=nextline]
\item[A1 --- the state.] An element's state is a normalised distribution of a unit over its poles, carried as amplitudes: a unit vector in $\R^N$ (\textbf{A1$\R$}) or in $\C^N$ (\textbf{A1$\C$}), where the $N$ poles are the outcome axes of its comparisons. \emph{Postulate}, except the choice of field: \textbf{complex amplitudes are an assumption}, and \S\ref{sec:geometry} shows the real primitives do not supply them.
\item[A2 --- the readable.] What is readable of a state is the set of its shares against every measurement---the transition shares $|\langle\varphi|\psi\rangle|^2$---and nothing else; frames and phases are not readable. \emph{Postulate.}
\item[A3 --- the relation.] A relation between two frames is a bijection of the readable states that preserves every transition share; the admissible relations form a \textbf{closed, connected group} of such bijections---closed under composition and inverse, containing the identity, and connected to it by continuous paths. \emph{Postulate}, with closedness and connectedness stated rather than inferred.
\item[A3b --- transitivity.] A relation may supply any coherence: the relations act transitively on the readable states. \emph{Postulate}---it is what ``supplied in neither pole'' means---but it is an axiom, and a reader who denies it keeps a closed connected subgroup and the theorems for that subgroup.
\item[A4 --- adjacency, the interface rule, locality.] Elements are joined by interfaces; the smallest closed paths of interfaces generate all closed paths; the cost of a configuration is the sum over smallest paths of a cost depending on that path alone. \textbf{The interface rule:} the drive pole an element receives along an interface \emph{is} the previous element's held pole carried across by the relation---an element has one held pole, its own, and a drive pole per arriving interface. When the poles are live, each element of a path reads the loop from its own pole and the path's cost is the mean of the readings. \emph{Assumption}: adjacency is the content of a lattice, and MINOS does not derive space.
\item[A5 --- the dynamics.] The dynamics is Hamiltonian, kinetic plus potential, additive over interfaces and smallest paths, with one coupling. \emph{Assumption}: the form is standard and not derived; the content of each term within the form is (\S\ref{sec:central}).
\item[A6 --- composition.] A composite of elements is an element over the union of their poles, with one normalisation. A constituent's state is the restriction of the composite's amplitude vector to the constituent's poles, subnormalised; its weight is the composite's datum. A constituent of zero weight is absent. \emph{Postulate} for restriction; \textbf{the gluing of independently specified constituents into one composite is not provided and is open} (\S\ref{sec:interp}).
\item[H1 --- the loss.] The readable loss of carrying a state round a closed path is one minus the returned share, and the loss of a path is that quantity averaged over the state with the invariant measure. \emph{Hypothesis}. Another monotone function of the share gives another class function, which this paper does not show to be equivalent. The average defines an averaged model (\S\ref{sec:central}), not a sector of a fuller one (\S\ref{sec:consequence}).
\end{description}

Seven axioms and one hypothesis. A1--A3b and A6 are the representation's own postulates: normalised amplitudes, transition shares as the readable, share-preserving bijections as relations, transitivity, composition by restriction. They are substantial mathematical commitments, and the representation \emph{proposes} them. A4 and A5 are the content of \emph{lattice field theory}. H1 is the choice of what to read. The honest form of the paper's claim is that the gauge structure lives entirely in the representation's postulates, and that lattice gauge theory is a severe test of what those postulates can do.

\section{Composing relations: the gauge geometry}\label{sec:geometry}

\subsection{The group}

\begin{theorem}\label{thm:group}
Under A1--A3b (A3 including closedness), A6 and the interface rule of A4:
\begin{enumerate}[label=(\roman*),leftmargin=2em]
\item (containment) the relations between frames of an $N$-pole element are induced by unitary maps of the amplitudes and act on the readable states as a closed connected subgroup $G\subseteq\PU(N)=\SU(N)/\mathbb{Z}_N$;
\item (the local freedom) every element's frame may be changed independently, so relations carry a local gauge freedom;
\item (the minimal real composite) a single two-pole element with real amplitudes has the full rotation group $\SO(2)$ of its circle; two such elements at an interface, composed under A6, have readable states forming the real projective plane $\RP^2$ with transition share $(n\cdot m)^2$, and relations generating $\SO(3)$ in its \textbf{vector} representation; the only other possibility for two real elements is one plane, $\SO(2)$. The qubit---complex amplitudes on $\CP^1$ with transition share $(1+n\cdot m)/2$---is \emph{not} obtained from two real elements;
\item (the complex branch) for an $N$-pole element with complex amplitudes, $G$ acts transitively on $\CP^{N-1}$ and is therefore $\PU(N)$ or, for even $N$, $\PSp(N/2)$, by Onishchik's classification \cite{onishchik};
\item (the real branch) for $N$ real poles, $G$ is a closed connected subgroup of $\SO(N)$ transitive on the sphere---$\SO(N)$ itself or one of the proper transitive subgroups of the Montgomery--Samelson--Borel list \cite{ms,borel}; for $N=3$ only $\SO(3)$; the group acting faithfully on $\RP^{N-1}$ is $G/(G\cap\{I,-I\})$---$\PSO(N)$ for even $N$ when $G=\SO(N)$.
\end{enumerate}
\end{theorem}

\begin{proof}[Proof (sketch; Appendix~\ref{app:proofs})]
(i) is Wigner's theorem \cite{wigner,bargmann} with antiunitaries excluded by continuity and phases quotiented by A2; (ii) is A2 applied to one element's frame; (iii) is a direct computation on the sphere of the composite's amplitudes under A6, with generation checked (check 10); (iv) and (v) are the classifications cited, applied because A3b makes the action transitive.
\end{proof}

Two consequences should be noted. On the real branch the non-Abelian structure is a property of a composite: one real element is $\SO(2)$, two are $\SO(3)$. On the complex branch a single two-pole element already has $\PU(2)$. And the complex amplitudes that give the qubit are nowhere derived. The representation therefore \textbf{forks}: real MINOS, with orthogonal groups on real projective spaces, and complex MINOS, with the groups of (iv). Nothing in A1--A6 selects between them. The fork is the source of the first result (\S\ref{sec:central}).

\subsection{The lattice construction, with the poles}\label{sec:construction}

Let the elements be the vertices $V$ of a graph and the interfaces its directed edges $E$. Each element $x$ carries, in its own frame, one \textbf{held pole} $c_x$---a unit amplitude vector read up to sign, a point of $\RP^{N-1}$ (real) or $\CP^{N-1}$ (complex)---which it hands across every departing interface. The relation on an edge $\ell=(x\to y)$ is an element $U_\ell\in G$, the map from $x$'s frame to $y$'s. The interface rule of A4 is the statement that the drive pole $y$ receives along $\ell$ is $x$'s held pole carried across:
\begin{equation}
d_{y,\ell}:=U_\ell c_x .
\end{equation}
The configuration space is therefore $G^E\times(\RP^{N-1})^V$ on the real branch and $G^E\times(\CP^{N-1})^V$ on the complex: the links together with the held poles. The drive poles are defined, not further variables, and this holds on every graph. A frame change $(g_x)_{x\in V}$ acts on amplitudes at $x$ by $g_x$, hence on poles by $c_x\mapsto g_xc_x$ and on relations by
\begin{equation}
U_{x\to y}\mapsto g_y\,U_{x\to y}\,g_x^{-1},
\end{equation}
the standard link transformation law, \emph{derived} from $U$ being the frame relation. The interface rule is preserved under it, so the rule is a transported label and constrains no link. The readables of an element are its \textbf{pass shares} $|\langle d_{y,\ell},c_y\rangle|^2$, one per arrival, and its \textbf{junction angles} $|\langle d_{y,\ell},d_{y,\ell'}\rangle|^2$, one per pair of arrivals. The pass share is the elemental angle of \S\ref{sec:comparison}, now seen to be a gauge invariant of the pair of elements. A vertex of higher degree adds readables, not constraints. The Hilbert space is $L^2\big(G^E\times(\RP^{N-1})^V\big)$ ($\CP^{N-1}$ on the complex branch), and the physical space is its $G^V$-invariant subspace. The generator at $x$ is the sum of the electric fields on the edges at $x$ plus the generator of rotations of $x$'s pole, and physical states are annihilated by it. This is \textbf{Gauss's law, with the pole's charge as its source}.

This is checked in finite dimensions, because it was contested. On the triangle (real, $N=3$), poles fixed in the frames independently of the links confine each link to a one-dimensional coset and leave no local freedom at all: three degrees of freedom, not a gauge theory. Poles transported as above leave every link free, with $\SO(3)^3$ acting freely, six degrees of freedom and six independent invariants. The fixed-pole case is recovered exactly as the transported one in the standard gauge with every pass share forced to zero (check 15). On four elements with two arrivals at one of them, the construction gives twenty variables, a free $\SO(3)^4$ orbit, eight degrees of freedom and eight invariants including the junction angle, against six under the alternative of forcing the arrivals to agree (check 16).

What the poles are: a director (real) or a ray (complex), whose frame-covariant content is the projector $Q_x=cc^{\mathsf T}$ or $\psi\psi^\dagger$, transforming as $Q_x\mapsto g_xQ_xg_x^{-1}$ in the \textbf{traceless symmetric} representation (real) or the \textbf{adjoint} (complex)---the very representations the loss will select in \S\ref{sec:central}. The projector is covariant, not invariant. The gauge invariants are traces of transported products of projectors. The simplest is the pass share $\tr(Q_yU_\ell Q_xU_\ell^{-1})$; the one the action reads is the loop's loss $1-\tr(Q_x\,WQ_xW^{-1})$, the failure of the pole's projector to return round the loop (check 16). The poles are projective matter; they are not a fundamental Higgs field, whose representation does not descend to $\PU(N)$.

\subsection{Composition, by identification or by relation}

A6 says a shared pole has one amplitude. Two systems that meet in a pole therefore compose \emph{by identification} only where they prescribe the same amplitude for it. Otherwise they compose \emph{by relation}: a junction whose readable is the overlap of the two arrivals. That overlap is one when they agree and zero when they are orthogonal, two systems that pass nothing to each other and do not collapse into one; between the two lies a share. This is why the lattice element has one held pole and a drive pole per arrival, and why a contested pole is not a contradiction. What A6 does not supply is the converse: a rule that glues independently specified comparisons, connected by arbitrary relations with arbitrary loop holonomies, into one composite state. Two constituents sharing several poles can prescribe incompatible amplitudes on them, and then no global vector exists. That gluing theorem is open, and is named in \S\ref{sec:interp} as the next task.

\section{The central result: under the loss rule, readable geometry determines the action's representation content}\label{sec:central}

\subsection{The two invariant measures}

Under A2, A3 and A3b the only measure on states that no frame can distinguish is the one invariant under all relations: with $G$ transitive and compact, the unique $G$-invariant probability measure. It differs between the branches. Under A1$\R$ it is the round measure on the real sphere of amplitudes, whose push-forward to the shares is the Fisher--Rao volume of the simplex (arcsine in the share for two poles). Under A1$\C$ it is the Fubini--Study volume on $\CP^{N-1}$, whose push-forward is uniform in the share. The two give different averages---for the return share under $\diag(i,-i)$, $1/2$ and $1/3$ respectively (check 12)---and the first drafts of this work ran them together.

\subsection{The complex branch: the adjoint action}

\begin{theorem}\label{thm:adjoint}
Under A1$\C$, A2--A4 and H1, for an $N$-pole element with relations $U\in\SU(N)$, the readable loss round a loop is
\begin{equation}
c(U)=\E_\psi\big[1-|\langle\psi|U|\psi\rangle|^2\big]=\frac{N-1}{N}\Big(1-\frac{\chi_{\adj}(U)}{N^2-1}\Big),
\end{equation}
the Wilson action in the adjoint representation.
\end{theorem}
\begin{proof}
$\E|\langle\psi|U|\psi\rangle|^2=(N+|\tr U|^2)/(N(N+1))$ by the second moment of the Fubini--Study measure, and $N\otimes\bar N=1\oplus\adj$ gives $|\tr U|^2=1+\chi_{\adj}$.
\end{proof}
(Check 8: $N=2$ to $5$ against Monte Carlo within error; the two forms equal to $10^{-16}$.) The readable sector is therefore centre-blind: the fundamental Wilson action is not a reading of the element's own shares but of a \emph{probe} that carries a frame of its own, and the string tension of such a probe is not a quantity of the readable sector.

\subsection{The real branch: the traceless symmetric action}

\begin{theorem2p}
Under A1$\R$, A2--A4, A6 and H1, for an element or composite of $N$ real poles with relations $R\in\SO(N)$ (or a transitive subgroup of Theorem~\ref{thm:group}(v)), the readable loss round a loop is
\begin{equation}
c_{\R}(R)=\E_n\big[1-(n\cdot Rn)^2\big]=\frac{2d_S}{N(N+2)}\Big(1-\frac{\chi_S(R)}{d_S}\Big),\qquad d_S=\frac{N(N+1)}{2}-1,\qquad \chi_S(R)=\frac{(\tr R)^2+\tr R^2}{2}-1,
\end{equation}
the Wilson action in the traceless symmetric representation $S$.
\end{theorem2p}
\begin{proof}
The fourth moment of the round measure gives $\E(n\cdot Rn)^2=((\tr R)^2+N+\tr R^2)/(N(N+2))$, and $\Sym^2(N)=1\oplus S$.
\end{proof}
For $N=3$ this is $\tfrac23(1-\chi_2/5)$, the \textbf{spin-two} action, and the adjoint form $\tfrac23(1-\chi_1/3)$ does not fit (check 11, four angles to the fourth decimal; check 13, $N=3$ to $6$ within Monte Carlo error). The adjoint never appears on the real branch, because a real share is a symmetric square.

\paragraph{The result.} The same operational construction---carry the readable state round the loop, read what returns, average over the frame no one can read---gives the adjoint character action when the amplitudes are complex and the traceless symmetric character action when they are real. The distinction between a representation and what is readable of it has a definite consequence: together with the loss rule, it decides which lattice action a system representation produces. Readability alone does not fix the action for every admissible loss; under H1 it fixes the representation content, and differently on the two branches. The character decompositions and moment formulae are known mathematics (the complex formula is the average-gate-fidelity integral) and we claim nothing for them; what is ours is the construction that selects them, and the fork it exposes.

\subsection{The kinetic term and the averaged Hamiltonian}\label{sec:averaged}

\begin{theorem}\label{thm:casimir}
The information distance between two relations---the distance between the share distributions they produce, averaged over the state---is bi-invariant on $G$, and the Laplace--Beltrami operator of the bi-invariant metric is the quadratic Casimir.
\end{theorem}
(Check 9: identical to six decimals after moving both relations by the same pair of group elements; small-displacement coefficient $1.00$ against the Killing form.)

\begin{theorem}[the averaged model]\label{thm:averaged}
Under A1--A5 and H1---the poles auxiliary, the loss of a path \emph{defined} as its average over the pole---the dynamics is
\begin{equation}
H=\frac{g^2}{2}\sum_{\text{interfaces}}\Delta_\ell+\frac{1}{g^2}\sum_{\text{smallest paths}}c(U_p),
\end{equation}
with $\Delta_\ell$ the Casimir on the interface's relation and $c$ the character loss of Theorem~\ref{thm:adjoint} or $2'$, on the states invariant under independent frame changes at every element: the Kogut--Susskind Hamiltonian \cite{ks} in the adjoint reading (traceless symmetric on the real branch), with Gauss's law the invariance condition.
\end{theorem}
\emph{This is a model defined by an averaging, and the theorem says so.} What the second result amounts to is therefore a reconstruction under explicit assumptions. Given A5's form, the content of each term is forced: the Casimir by the metric, the character action by the loss. The form itself is not derived, and neither is the claim that nature's dynamics is generated by this Laplacian rather than by a function of it. (Locality, gauge invariance and continuity in the metric do not select it: $\sum(-\Delta_\ell)^2+\sum c(U_p)$ satisfies all three.)

The two readings meet at weak coupling: the fundamental and adjoint actions have the same leading local continuum term after coupling normalisation, with couplings in the ratio $2N/(N+1)$ (check 4), which is a statement about the local term and not about the global structure of $\SU(N)$ against $\PU(N)$.

\section{A nontrivial consequence: the projective-pole model and its induced interaction}\label{sec:consequence}

\subsection{The poles kept live}

Averaging the loss over the pole and restricting the dynamics to pole-independent states are different operations. With $P$ the projector on states independent of the poles, $PVP=c(U)P$---the character action---but $(I-P)VP\neq0$: on one $\SO(3)$ plaquette with one pole, for the quarter turn, $PVP=0.800000=\tfrac23(1-\chi_2/5)$ while $\|(I-P)VP\|=0.2667$, the standard deviation of the loss over the pole (check 17). A pole-independent state leaves the subspace at once. \textbf{There is no pole-neutral sector of the full dynamics.} The pure theory of Theorem~\ref{thm:averaged} is the averaged model, or---as follows---the leading term of a limit.

\subsection{The reduction, and the two-plaquette kernel}

\begin{theorem4p}
Let the poles be dynamical with their own kinetic term---by Theorem~\ref{thm:casimir} on the pole's state space, the Laplacian of its invariant metric, with coefficient $\kappa$---so that the full Hamiltonian is
\begin{equation}
H_{\mathrm{full}}=\frac{g^2}{2}\sum_\ell\Delta_\ell+\kappa\sum_x\Delta_{\mathrm{pole},x}+\frac{1}{g^2}\sum_pV_p,\qquad V_p=\frac{1}{|p|}\sum_{x\in p}V\big(W_p^{(x)},n_x\big),
\end{equation}
with $V_p$ the loss of path $p$ read from each of its $|p|$ elements and averaged (A4). Then, as a formal low-energy asymptotic expansion for a fixed finite graph in a window of bounded link energy,
\begin{equation}
E_0(U;\kappa)=\frac{1}{g^2}\sum_pc(U_p)-\frac{1}{\kappa g^4}\sum_x\sum_{p,q\ni x}\frac{K\big(V_p^{(x)},V_q^{(x)}\big)}{|p|\,|q|}+O(\kappa^{-2}),\qquad
K(V_1,V_2)=\sum_{m\neq0}\frac{\langle0|V_1|m\rangle\langle m|V_2|0\rangle}{\lambda_m},
\end{equation}
the inner sum over ordered pairs of smallest paths through $x$ and $\lambda_m$ the pole Laplacian's eigenvalues. The leading term is the character action. The next is a sum of one-plaquette corrections ($p=q$), which are class functions, and of \textbf{two-plaquette interactions} ($p\neq q$) wherever two paths are read through one pole. Each two-plaquette term is a function of the pair of holonomies based at the shared element; it is invariant under their simultaneous conjugation and not determined by their separate classes.
\end{theorem4p}

\paragraph{The kernel in closed form (real branch, $N=3$).} For a rotation through $\varphi$ about the axis $a$, $n\cdot Rn=t+(1-t)(n\cdot a)^2$ with $t=\cos\varphi$, so the loss has harmonics of degree $0$, $2$ and $4$ only, with zonal coefficients
\begin{equation}
A_2(t)=-\tfrac{4}{21}(1-t)(4t+3),\qquad A_4(t)=-\tfrac{8}{35}(1-t)^2,
\end{equation}
and by the addition theorem
\begin{equation}\label{eq:kernel}
K(R_1,R_2)=\frac{A_2(t_1)A_2(t_2)}{30}\,P_2(a\cdot b)+\frac{A_4(t_1)A_4(t_2)}{180}\,P_4(a\cdot b).
\end{equation}
For two quarter turns, $K=\tfrac{8}{735}P_2(a\cdot b)+\tfrac{16}{55125}P_4(a\cdot b)$: $+0.0111746$ for parallel axes, $-0.0053333$ for perpendicular; a quarter turn beside a half turn gives $-16/2625$ even for parallel axes. Under A4's average a plaquette of $L$ elements has its own second-order coefficient $K_1/L$, $0.00279365$ for a quarter-turn square.

\paragraph{What is fixed.} The two-plaquette interaction has no coefficient of its own. Its size relative to the one-plaquette correction, its sign, and its dependence on the two angles and the relative axis are all fixed by the pole being projective matter in the vector representation with H1's quartic loss. The free numbers of the model are $g$ and $\kappa$. The stiffness $\kappa$ is a second number the axioms do not fix. This is the third result: once the model is specified, these quantities are determined. Integrating out stiff matter to induce gauge interactions is an established subject \cite{km,ps}; we have neither found nor excluded an exact prior match for this kernel, and we claim only that the construction fixes it.

\paragraph{What the model is, and is not.} $H_{\mathrm{full}}$ gives the poles a kinetic term and the loop losses, and nothing else. With every holonomy trivial the potential vanishes whatever the poles' arrangement, so the projective field has no spatial stiffness of its own. A hopping term $\tr(Q_yU_\ell Q_xU_\ell^{-1})$, the gauged $\CP^{N-1}$ model's \cite{cpn}, is gauge-invariant and available. A4 and A5 put the potential on closed paths only, which is why it is absent; nothing in the axioms derives such a term and nothing forbids it. The stated theory is a specialised member of the family ``gauge theory coupled to projective matter''. Making the poles stiff is not a continuum limit: $\kappa\to\infty$ gives the pure lattice theory on the graph, and nothing here concerns four dimensions or $a\to0$.

\section{Verification}\label{sec:verification}

Every statement above with a number in it is produced by the one accompanying script, seeded, in about fifteen minutes with numpy and scipy, and written to \texttt{results.txt} and \texttt{results.json} beside it. The checks that carry the argument:

\begin{center}\small
\begin{tabular}{@{}>{\raggedright\arraybackslash}p{1.4cm}>{\raggedright\arraybackslash}p{7.6cm}>{\raggedright\arraybackslash}p{6.2cm}@{}}
\toprule
check & what it establishes & result\\
\midrule
8 & Theorem~\ref{thm:adjoint}'s second moment and the adjoint form, $N=2$--$5$ & Monte Carlo within error; forms equal to $10^{-16}$\\
11, 13 & Theorem $2'$: the spin-two action for $N=3$; the traceless symmetric action for $N=3$--$6$ & to the fourth decimal; within error\\
12 & the two invariant measures told apart on $\diag(i,-i)$ & $1/3$ and $1/2$\\
9 & Theorem~\ref{thm:casimir}: bi-invariance and the Casimir coefficient & six decimals; $1.00$\\
15 & the interface rule against the free link on the triangle & fixed poles 3 d.o.f.; transported 6; coset $=$ zero-pass-share sector\\
16 & the general graph with two arrivals at one element & 8 d.o.f., 8 invariants; the constrained alternative 6\\
17 & no pole-neutral sector; the one-pole stiffness expansion & $PVP=0.8$; $\|(I-P)VP\|=0.2667$; $\kappa(c-E_0)\to0.01118$ against $0.01117$\\
18, 19(a) & the two-plaquette kernel, quarter turns at five angles and arbitrary angles & to $10^{-16}$; $-16/2625$ recovered\\
19(b) & the complete A4-averaged fixed-link Hamiltonian on two squares sharing an element, every pole kept & $K_1/4.00$; residual$\times\kappa^2=-0.00004$ at every $\kappa$ from 10 to 160\\
20, 21 & the reduction with the link dynamical: one plaquette's link-and-pole Hamiltonian diagonalised directly against the effective operator, across truncations to 18\,876 states, in the gauge-invariant sector, four lowest levels & residual$\times\kappa^2$ steady at $0.0022$--$0.0028$ across $\kappa=5$--$80$\\
\bottomrule
\end{tabular}
\end{center}

The strongest of these are 19(b) and 21: the first tests the whole finite-graph Hamiltonian rather than its components, the second lets the link move quantum mechanically and compares spectra rather than one energy. Both are numerical evidence on fixed graphs, not operator theorems, and are described so. Two instrument faults found during the checks are kept in the results with their corrected forms. One was an effective operator first compared against the exact kernel function rather than one formed in the same truncation; it produced a $1/\kappa$ residual at small $g$ that was the truncation's. The other was a generator convention that failed to commute with the potential; it was found by testing the commutator and replaced by the one that commutes to $10^{-13}$. Appendix~\ref{app:loop} gives the adjunct results on the smallest closed loop with the same standard.

\section{Interpretation and limits}\label{sec:interp}

\subsection{Three kinds of result}

\begin{center}\small
\begin{tabular}{@{}>{\raggedright\arraybackslash}p{3.6cm}>{\raggedright\arraybackslash}p{11.6cm}@{}}
\toprule
category & what it is here\\
\midrule
\textbf{established mathematics} & Wigner's theorem and the transitive-group classifications; character decomposition and invariant integration; the Casimir as the bi-invariant Laplacian; degenerate perturbation theory and the Schur complement\\
\addlinespace
\textbf{the representation's contribution} & an operational organisation of these tools around a primitive comparison, a rule of what is readable, and relations supplied between frames---under which the group, the local freedom, Gauss's law, the measure, the actions and the kinetic term are produced in sequence, with the real/complex fork exposed\\
\addlinespace
\textbf{demonstrated usefulness} & the construction yields coherent gauge-invariant models on any graph; what is readable selects representation-dependent actions; with the poles live, it fixes a particular induced interaction whose form is not a free choice; and every step was made precise enough to be checked, and was\\
\bottomrule
\end{tabular}
\end{center}

\subsection{What is derived, reconstructed, and model-dependent}

\emph{Derived} from the representation's postulates: the link variable is a group element and which group (Theorem~\ref{thm:group}); the local gauge freedom and the transformation law; Gauss's law with the pole charge; the invariant measure; the character actions under H1 (Theorems~\ref{thm:adjoint}, $2'$); the Casimir as the natural Laplacian (Theorem~\ref{thm:casimir}). \emph{Reconstructed} under assumptions: the Hamiltonian (Theorem~\ref{thm:averaged} needs A5's form and H1's average). \emph{Consequences of the specified model}: Theorem $4'$ and its kernel. \emph{Hypotheses throughout}: complex amplitudes; transitivity; adjacency; the kinetic-plus-potential form; the loss functional; the poles' stiffness.

\subsection{What remains open, and is not claimed}

\begin{itemize}[leftmargin=1.5em]
\item \textbf{The composition law.} A6 restricts a given composite; it does not glue independently specified comparisons with arbitrary holonomies into one. Defining that operation, and finding which constraints it imposes, is the next mathematical task; it is also the point at which the representation's own content, rather than its organisation of known content, would be tested.
\item \textbf{Complex amplitudes} remain a hypothesis; the minimal real composite gives $\SO(3)$ in its vector representation on $\RP^2$ and a spin-two action, not the qubit.
\item \textbf{Whether the fixed interaction is a restriction} competing gauge--matter models lack, or a rediscovery within them, awaits a literature comparison we have not made.
\item \textbf{The continuum theory and the mass gap} are untouched. Nothing here bears on four dimensions, on $a\to0$, or on the spectrum of the continuum theory, and we make no claim about them.
\end{itemize}

\subsection{Conclusion}

The results establish a qualified but substantive first test of the representation as an operational account of systems. With its mathematical postulates and its domain assumptions stated explicitly, the framework provides a coherent construction of gauge-covariant relations, representation-dependent readable observables and the corresponding lattice actions. In a specified dynamical projective-pole model it additionally determines higher-order interactions that admit direct analytical and numerical verification.

These results do not establish a new foundation for Yang--Mills theory or a new physical prediction. They demonstrate something more limited and independently useful. A representation devised for general systems can be applied to an established mathematical theory, made precise, and used to construct and analyse nontrivial structures without merely renaming the theory's existing objects. In the other direction, that theory's structures can be understood as consequences of a few statements about systems, which is the comprehensional purpose the paper set itself.

The broader validity of the representation as an account of systems in general remains open. So does the question that would test its own content rather than its organisation of known content: whether its compositional principles impose restrictions beyond those available in conventional gauge theory. This paper is the first case study in which the representation is stated explicitly enough for that question to be put.

\appendix

\section{Proofs}\label{app:proofs}

\paragraph{Theorem~\ref{thm:group}.} (i) By A2 the readable is the set of transition shares; by A3 a relation preserves all of them; Wigner's theorem states that a bijection of projective space preserving all transition shares is induced by a unitary or antiunitary map, unique up to phase; antiunitaries are not connected to the identity and are excluded by A3's continuity; a phase induces the same map on readable states, so relations act through $\PU(N)$, and the admissible set is a closed (A3) connected (continuous composition from the identity) subgroup. (ii) A change of frame at $x$ applied to every relation touching $x$ changes nothing readable, so $U_\ell\mapsto g_yU_\ell g_x^{-1}$ for independent $g$ is a symmetry. (iii) A real two-pole element has states on the unit circle read up to sign, $\RP^1$, with readable-preserving connected maps $\SO(2)$, transitive by A3b. Two such elements $A$ (poles $x,y$) and $B$ (poles $y,z$) at an interface compose under A6 into an element over $(x,y,z)$ with one normalisation: amplitudes a unit vector $n\in S^2$, read up to sign, $\RP^2$; $A$'s state is the restriction $(n_x,n_y)$ with weight $n_x^2+n_y^2$, $B$'s $(n_y,n_z)$---defined everywhere, absent exactly at the measure-zero points where the weight vanishes. $A$'s relations rotate the $(x,y)$ plane, $B$'s the $(y,z)$ plane; rotations about two distinct axes generate $\SO(3)$ (check 10); the transition share between readable states $\pm n$, $\pm m$ is $(n\cdot m)^2$, the vector representation. $\RP^2$ and $\CP^1$ are not homeomorphic and their transition shares differ, so the qubit is not obtained. (iv) $G\subseteq\PU(N)$ acts transitively on $\CP^{N-1}$ by A3b; the compact connected groups transitive on $\CP^{N-1}$ are $\SU(N)$ and, for even $N$, $\Sp(N/2)$, up to the kernel \cite{onishchik}. (v) On the real branch the connected readable-preserving maps are $\SO(N)$; A3b makes the action transitive on $\RP^{N-1}$, hence on $S^{N-1}$; the transitive compact connected groups on spheres are classified \cite{ms,borel}; $-I$ acts trivially on $\RP^{N-1}$.\qed

\paragraph{Theorem~\ref{thm:adjoint}.} For $\psi$ Fubini--Study distributed, $\E[\psi_i\bar\psi_j\psi_k\bar\psi_l]=(\delta_{ij}\delta_{kl}+\delta_{il}\delta_{kj})/(N(N+1))$; contracting with $U_{ji}\bar U_{lk}$ gives $\E|\langle\psi|U|\psi\rangle|^2=(N+|\tr U|^2)/(N(N+1))$; $N\otimes\bar N=1\oplus\adj$ gives $|\tr U|^2=1+\chi_{\adj}(U)$.\qed

\paragraph{Theorem $2'$.} $\E[n_in_jn_kn_l]=(\delta_{ij}\delta_{kl}+\delta_{ik}\delta_{jl}+\delta_{il}\delta_{jk})/(N(N+2))$ on $S^{N-1}$; so $\E(n\cdot Rn)^2=((\tr R)^2+N+\tr R^2)/(N(N+2))$; $\Sym^2(N)=1\oplus S$ with character $((\tr R)^2+\tr R^2)/2=1+\chi_S$; hence $1-\E(n\cdot Rn)^2=(2/(N(N+2)))(d_S-\chi_S)$. For a proper transitive subgroup the round measure is still the invariant one and the formula holds with $\chi_S$ restricted.\qed

\paragraph{Theorem~\ref{thm:casimir}.} The distance between the share distributions of two relations $U,U'$, averaged over the state with the invariant measure, is a function of $U'U^{-1}$ invariant under conjugation, hence bi-invariant; the bi-invariant metric on a compact simple group is the Killing form up to scale, and its Laplace--Beltrami operator is the quadratic Casimir \cite{helgason}.\qed

\paragraph{Theorem~\ref{thm:averaged}.} Kinetic term: Theorem~\ref{thm:casimir} with A5's form. Potential: Theorem~\ref{thm:adjoint} or $2'$ with A4, a function of the links alone under H1's average. Invariance: by A2 two states of the relations that differ by a frame change at an element are one readable state; the generators---the electric fields summed at an element---annihilate physical states.\qed

\paragraph{Theorem $4'$.} At fixed links the poles decouple, and each pole's Hamiltonian is $\kappa\Delta+g^{-2}\sum_{p\ni x}V_p^{(x)}/|p|$. Degenerate perturbation theory about the uniform ground state (gap $\kappa\lambda_1$) gives, at second order, $-(\kappa g^4)^{-1}\sum_{m\neq0}\big|\langle m|\sum_pV_p^{(x)}/|p|\,|0\rangle\big|^2/\lambda_m$, whose expansion of the square gives the double sum; poles at different elements are independent, so only paths through a common element couple; invariance under simultaneous conjugation follows from the invariance of the measure and the Laplacian. With the links dynamical the same expansion is the Schur complement of the pole-excited block, whose resolvent is $1/(\kappa\lambda_m)$ plus terms carrying the link Laplacian, of relative order $g^2/(\kappa\lambda_m)$, entering at $O(\kappa^{-2})$ within a window of bounded link energy; checks 20 and 21 test this and the operator statement is not claimed beyond that window. The closed-form kernel: $V_{\varphi,a}(n)=1-[t+(1-t)(n\cdot a)^2]^2$ expanded in Legendre polynomials of $n\cdot a$ gives the zonal coefficients $A_2$, $A_4$; the addition theorem $\sum_mY_{\ell m}(a)\bar Y_{\ell m}(b)=(2\ell+1)P_\ell(a\cdot b)/4\pi$ and $\lambda_\ell=\ell(\ell+1)$ give $K$.\qed

\section{The smallest closed loop (adjunct results)}\label{app:loop}

For $\SU(2)$ and one loop, $H=\lambda C_2+v(U)/\lambda$ on class functions, with $v$ the amplitude (fundamental) or the share (adjoint) cost matched at weak coupling (check 6). In the share reading the potential never populates half-integer $j$: the vacuum sector of the readable reading contains only whole carriers, an exact selection rule at every coupling. At strong coupling the ratio of the two readings' gaps tends to $C_2(1)/C_2(\tfrac12)=8/3=C_A/C_F$---the first excitation of a loop is the fundamental carrier under the amplitude reading and the adjoint carrier under the share reading, at strong coupling---and at weak coupling the ratio tends to one, as the common leading local term requires. The lattice values of the lightest glueball mass in units of the string tension across $\SU(N)$ \cite{lt,teper} vary with $N$ in the direction, and roughly by the amount, that $\sqrt{C_A/C_F}$ would give; in this representation that would read as the carrier of the lowest excitation being the adjoint one. We record the comparison as a reading: the proportionality constant is not derived, and the scaling is not claimed as a result here.

Prepared as the ground state at coupling $2\lambda$ and evolved at coupling $\lambda$, the loop is neutral at every instant---neutrality is preserved by the gauge-invariant evolution---and the expectation of its held share oscillates about a non-zero mean. Its dominant frequency is the gap of $H$ to four decimals at $\lambda=0.5,1,2$; the exact spectral content (check 7) is that the prepared state's weight in the ground and first connected excited state is $0.99998$, $1.000000$, $1.000000$ and the gap frequency carries the same share of the alternating power, the remainder sitting at the other transition frequencies. The two-level description holds to those figures and no better, and an exact single-frequency statement would need a selection rule we do not have.

\begin{thebibliography}{12}
\bibitem{wigner} E.~P. Wigner, \emph{Gruppentheorie und ihre Anwendung auf die Quantenmechanik der Atomspektren}, Vieweg (1931).
\bibitem{bargmann} V.~Bargmann, Note on Wigner's theorem on symmetry operations, \emph{J. Math. Phys.} \textbf{5} (1964) 862.
\bibitem{helgason} S.~Helgason, \emph{Groups and Geometric Analysis}, Academic Press (1984).
\bibitem{ks} J.~Kogut and L.~Susskind, Hamiltonian formulation of Wilson's lattice gauge theories, \emph{Phys. Rev. D} \textbf{11} (1975) 395.
\bibitem{bc} G.~Bhanot and M.~Creutz, Variant actions and phase structure in lattice gauge theory, \emph{Phys. Rev. D} \textbf{24} (1981) 3212.
\bibitem{onishchik} A.~L. Onishchik, Transitive compact transformation groups, \emph{Mat. Sb.} \textbf{60} (1963) 447; \emph{Amer. Math. Soc. Transl.} \textbf{55} (1966) 153.
\bibitem{fs} E.~Fradkin and S.~H. Shenker, Phase diagrams of lattice gauge theories with Higgs fields, \emph{Phys. Rev. D} \textbf{19} (1979) 3682---from which the projective poles are distinguished.
\bibitem{ms} D.~Montgomery and H.~Samelson, Transformation groups of spheres, \emph{Ann. of Math.} \textbf{44} (1943) 454.
\bibitem{borel} A.~Borel, Some remarks about Lie groups transitive on spheres and tori, \emph{Bull. Amer. Math. Soc.} \textbf{55} (1949) 580.
\bibitem{cpn} A.~D'Adda, M.~L\"uscher and P.~Di~Vecchia, A $1/n$ expandable series of non-linear $\sigma$ models with instantons, \emph{Nucl. Phys. B} \textbf{146} (1978) 63.
\bibitem{km} V.~A. Kazakov and A.~A. Migdal, Induced QCD at large $N$, \emph{Nucl. Phys. B} \textbf{397} (1993) 214.
\bibitem{ps} A.~Pap and P.~Suranyi, Gauge theories induced by bosons, hep-lat/9608112.
\bibitem{lt} B.~Lucini and M.~Teper, SU($N$) gauge theories in four dimensions: exploring the approach to $N=\infty$, \emph{JHEP} \textbf{06} (2001) 050.
\bibitem{teper} M.~Teper, Glueball masses and other physical properties of SU($N$) gauge theories in $D=3+1$, hep-th/9812187.
\end{thebibliography}

\end{document}
